% ===================================================================
%  ਸਾਰਣੀ ਨੰ. 6 — ਘਰ ਦੇ ਅੰਦਰ। — ਸਾਜ਼-ਸਾਮਾਨ। — ਖਾਣਾ। — ਚਾਨਣ।
%  Traduction 2026 d'apres texto/io, controlee sur texto/fr, texto/en
%  et texto/hi. Consigne : voir texto/pa/10-sarni-01.tex.
%
%  Deux notes, toutes deux marquees « (*) », et deux appels : celui de
%  « babo » au bloc c1-12-1, celui de « lambrequino » au bloc c2-01-1.
%  L'ordre les departage, comme dans les dix autres colonnes.
%
%  LA FEMME DE CHAMBRE N'A PAS DE GROS PLAN, et c'est normal. L'ido du
%  bloc c1-06-1 porte « la chambristino (6) », le francais « la femme
%  de chambre (16) » : c'est l'ido qui se trompe, son propre numero 6
%  etant le savon du bloc c1-02-1. C'est un des trois ecarts declares
%  dans outils/renvoji.py.
% ===================================================================

%%K t06-apar-1 apar
\VUsaut{6.0mm}
\VUtitre{15.0pt}{60}{ਸਾਰਣੀ ਨੰ. 6}
\VUsaut{1.4mm}
\VUtitre{11.4pt}{0}{ਘਰ ਦੇ ਅੰਦਰ, ਸਾਜ਼-ਸਾਮਾਨ, ਖਾਣਾ,}
\VUsaut{0.4mm}
\VUtitre{11.4pt}{0}{ਚਾਨਣ।}
\VUsaut{2.2mm}

%%K t06-apar-2 apar
ਘਰ ਬਣ ਕੇ ਤਿਆਰ ਹੈ। ਜਾਨ ਦੇ ਚਾਚਾ ਆਪਣੇ ਨਵੇਂ ਘਰ ਵਿਚ ਆ ਵੱਸੇ ਹਨ, ਜਿਸ ਦੀ ਬਣਤਰ ਅਸੀਂ
ਜਾਣਦੇ ਹੀ ਹਾਂ। ਹੁਣ ਵੇਖੀਏ ਕਿ ਉਨ੍ਹਾਂ ਨੇ ਉਸ ਨੂੰ ਕਿਵੇਂ ਸਜਾਇਆ ਹੈ। ਪਹਿਲਾਂ ਅਸੀਂ
ਵੇਖਾਂਗੇ:

%%K t06-tit-1 sub
\VUpk{10.2pt}{40}{ਸੌਣ ਦਾ ਕਮਰਾ।}
\VUsaut{3.0mm}

%%K t06-c1-01-1 p
1. --- ਅਸੀਂ ਕੁਝ ਬੇਵੇਲੇ ਆ ਗਏ ਹਾਂ, ਕਿਉਂਕਿ ਨੌਕਰਾਣੀ \VUgras{ਕਮਰੇ} ਦੀ ਸਫ਼ਾਈ ਵਿਚ
ਲੱਗੀ ਹੈ।

%%K t06-c1-02-1 p
2. --- \VUgras{ਹੱਥ-ਮੂੰਹ ਧੋਣ ਦੀ ਚੌਕੀ} \textsuperscript{(1)} ਉੱਤੇ ਇੱਧਰ ਉੱਧਰ ਉਹ
ਚੀਜ਼ਾਂ ਪਈਆਂ ਹਨ ਜਿਨ੍ਹਾਂ ਨਾਲ ਨ੍ਹਾਇਆ, ਕੰਘੀ ਕੀਤੀ ਤੇ ਤਿਆਰ ਹੋਇਆ ਜਾਂਦਾ ਹੈ। ਉਹ ਹਨ
\VUgras{ਚਿਲਮਚੀ} \textsuperscript{(2)} ਤੇ \VUgras{ਪਾਣੀ ਦਾ ਜੱਗ} \textsuperscript{(3)},
\VUgras{ਵਾਲਾਂ ਦਾ ਬੁਰਸ਼} \textsuperscript{(4)}, \VUgras{ਕੰਘੀ} \textsuperscript{(5)},
\VUgras{ਸਾਬਣ} \textsuperscript{(6)} ਤੇ \VUgras{ਸਪੰਜ} \textsuperscript{(7)}, ਤੇ ਅੰਤ
ਵਿਚ ਕੁਰਲੀ ਕਰਨ ਦੇ ਪਾਣੀ ਦੀ ਇਕ ਨਿੱਕੀ \VUgras{ਸ਼ੀਸ਼ੀ} \textsuperscript{(8)}।

%%K t06-c1-03-1 p
3. --- ਭੁੰਜੇ, ਚੌਕੀ ਦੇ ਸਾਹਮਣੇ, ਇਹ ਰਹੀ ਗੰਦੇ ਪਾਣੀ ਦੀ \VUgras{ਬਾਲਟੀ}
\textsuperscript{(9)}।

%%K t06-c1-04-1 p
4. --- \VUgras{ਬਾਹਰਲੇ ਕਮਰੇ} \textsuperscript{(10)} ਵਿਚ ਸਾਨੂੰ ਕੁਝ
\VUgras{ਖੂੰਟੀਆਂ} \textsuperscript{(13)} ਵਿਖਾਈ ਦਿੰਦੀਆਂ ਹਨ ਤੇ ਦੋ
\VUgras{ਛਤਰੀਆਂ} \textsuperscript{(11)}, ਜੋ \VUgras{ਛਤਰੀ-ਦਾਨ}
\textsuperscript{(12)} ਵਿਚ ਹਨ।

%%K t06-c1-05-1 p
5. --- ਬੂਹੇ ਦੇ ਦੂਜੇ ਪਾਸੇ \VUgras{ਸ਼ੀਸ਼ੇ ਵਾਲੀ ਅਲਮਾਰੀ} \textsuperscript{(14)} ਖੜੀ
ਹੈ, ਜਿਸ ਵਿਚ \VUgras{ਕੱਪੜੇ} \textsuperscript{(15)} ਰੱਖੇ ਰਹਿੰਦੇ ਹਨ।

%%K t06-c1-06-1 p
6. --- \VUgras{ਨੌਕਰਾਣੀ} \textsuperscript{(16)} \VUgras{ਪਲੰਘ}
\textsuperscript{(17)} ਲਾ ਰਹੀ ਹੈ, ਇਸੇ ਬਹਾਨੇ ਉਸ ਦੇ ਵੱਖੋ-ਵੱਖ ਹਿੱਸੇ ਵੇਖ ਲਈਏ।
\VUgras{ਮੰਜਾ} \textsuperscript{(20)} ਉੱਤੇ ਇਕ ਚੰਗਾ \VUgras{ਕਮਾਨੀ ਵਾਲਾ ਗੱਦਾ}
\textsuperscript{(21)} ਹੈ, ਜਿਸ ਉੱਤੇ ਜਸਟੀਨ ਹੁਣੇ ਉੱਨ ਦਾ \VUgras{ਗੱਦਾ}
\textsuperscript{(22)} ਵਿਛਾਏਗੀ। ਫਿਰ ਉਹ ਕੁਰਸੀਆਂ ਤੋਂ ਦੋ \VUgras{ਚਾਦਰਾਂ}
\textsuperscript{(23)} ਤੇ \VUgras{ਕੰਬਲ} \textsuperscript{(24)} ਚੁੱਕੇਗੀ। ਫਿਰ ਉਹ ਪਲੰਘ
ਦੇ ਸਿਰਹਾਣੇ \VUgras{ਲੰਮਾ ਸਿਰਹਾਣਾ} \textsuperscript{(25)} ਤੇ \VUgras{ਸਿਰਹਾਣਾ}
\textsuperscript{(26)} ਰੱਖੇਗੀ, ਜੋ ਹਾਲੇ \VUgras{ਰਜਾਈ} \textsuperscript{(27)} ਨਾਲ
\VUgras{ਪਲੰਘ ਦੇ ਅੱਗੇ ਦੀ ਦਰੀ} \textsuperscript{(32)} ਉੱਤੇ ਪਏ ਹਨ। ਉਹ
\VUgras{ਪਰਦਿਆਂ} \textsuperscript{(19)} ਤੇ \VUgras{ਛੱਤਰ} \textsuperscript{(18)} ਉੱਤੇ
ਖੰਭਾਂ ਦਾ ਝਾੜਨ ਫੇਰਦੀ ਹੈ, ਤੇ ਅੱਧਾ ਕੰਮ ਹੋ ਜਾਂਦਾ ਹੈ।

%%K t06-c1-07-1 p
7. --- ਫਿਰ ਉਸ ਨੂੰ \VUgras{ਪਲੰਘ ਕੋਲ ਦੀ ਮੇਜ਼} \textsuperscript{(28)} ਥੋੜੀ ਸੰਵਾਰਨੀ
ਪਵੇਗੀ, \VUgras{ਅਲਾਰਮ ਘੜੀ} \textsuperscript{(29)} ਵਿਚ ਚਾਬੀ ਭਰਨੀ ਪਵੇਗੀ,
\VUgras{ਰਾਤ ਦੀ ਬੱਤੀ} \textsuperscript{(30)} ਤਿਆਰ ਕਰਨੀ ਪਵੇਗੀ, ਤੇ
\VUgras{ਮੋਮਬੱਤੀ} \textsuperscript{(31)} ਹਟਾ ਕੇ ਸ਼ਮ੍ਹਾਦਾਨ ਸਾਫ਼ ਕਰਨਾ ਪਵੇਗਾ।

%%K t06-tit-2 sub
\VUsaut{5.0mm}
\VUpk{10.2pt}{40}{ਸਿਲਾਈ ਦੀ ਟੋਕਰੀ ਤੇ ਅੰਗੀਠੀ ਦਾ ਸਾਮਾਨ।}
\VUsaut{3.0mm}

%%K t06-c1-08-1 p
8. --- \VUgras{ਸਿਲਾਈ ਦੀ ਟੋਕਰੀ} \textsuperscript{(34)}, ਜੋ \VUgras{ਤਿਪਾਈ}
\textsuperscript{(33)} ਕੋਲ ਰੱਖੀ ਹੈ, ਉਸ ਨੂੰ ਵੀ ਸੰਵਾਰਨ ਦੀ ਬੜੀ ਲੋੜ ਹੈ। ਉਸ ਨੂੰ
\VUgras{ਕਢਾਈ} \textsuperscript{(35)} ਤਹਿ ਕਰਨੀ ਪਵੇਗੀ, \VUgras{ਸਿਲਾਈ ਦੀ ਮੇਜ਼}
\textsuperscript{(38)} ਵਿਚ \VUgras{ਅੰਗੂਠਾ}, \VUgras{ਧਾਗਾ} \textsuperscript{(36)},
\VUgras{ਸੂਈਆਂ} \textsuperscript{(37)} ਤੇ \VUgras{ਕੈਂਚੀ} \textsuperscript{(39)}
ਰੱਖਣੀਆਂ ਪੈਣਗੀਆਂ, ਤੇ ਮੇਜ਼ ਦਾ ਢੱਕਣ \VUgras{ਚਾਬੀ} \textsuperscript{(40)} ਨਾਲ ਬੰਦ ਕਰਨਾ
ਪਵੇਗਾ।

%%K t06-c1-09-1 p
9. --- ਵਿਚਕਾਰ ਰੱਖੀ \VUgras{ਤਿਪਾਈ-ਮੇਜ਼} \textsuperscript{(41)} ਉੱਤੇ ਜਸਟੀਨ
\VUgras{ਸੁਰਾਹੀ} \textsuperscript{(42)} ਦਾ ਪਾਣੀ ਬਦਲੇਗੀ। ਉਹ \VUgras{ਖੰਡ-ਦਾਨ}
\textsuperscript{(43)} ਵਿਚ \VUgras{ਖੰਡ} ਪਾਵੇਗੀ, \VUgras{ਖੰਡ ਦੀ ਚਿਮਟੀ}
\textsuperscript{(44)} ਸਾਫ਼ ਕਰੇਗੀ ਤੇ \VUgras{ਕੱਪੜਿਆਂ ਦਾ ਬੁਰਸ਼}
\textsuperscript{(46)} ਹਟਾ ਦੇਵੇਗੀ।

%%K t06-c1-10-1 p
10. --- ਕਮਰੇ ਦਾ ਸੱਜਾ ਪਾਸਾ ਉਸ ਨੂੰ ਘੱਟ ਕੰਮ ਦੇਵੇਗਾ, ਕਿਉਂਕਿ \VUgras{ਦਿਨ ਦਾ ਪਲੰਘ}
\textsuperscript{(47)} ਤੇ ਉਸ ਦੀ \VUgras{ਗੱਦੀ} \textsuperscript{(48)} ਆਪਣੀ ਥਾਂ ਉੱਤੇ
ਹਨ, ਤੇ ਕਿਉਂਕਿ ਮੌਸਮ ਠੰਢਾ ਨਹੀਂ, \VUgras{ਅੰਗੀਠੀ} \textsuperscript{(49)} ਦੇ ਸਾਮਾਨ
ਵਿਚ ਕੁਝ ਹਿਲਾਇਆ ਨਹੀਂ ਗਿਆ। \VUgras{ਕਹੀ} \textsuperscript{(50)}, \VUgras{ਚਿਮਟਾ}
\textsuperscript{(51)}, \VUgras{ਲੱਕੜ ਦੇ ਟੇਕ} \textsuperscript{(55)} ਤੇ
\VUgras{ਸੁਆਹ ਦਾ ਪਰਦਾ} \textsuperscript{(56)} ਸਾਫ਼ ਹਨ; \VUgras{ਅੱਗ ਦਾ ਪਰਦਾ}
\textsuperscript{(54)} ਹਿਲਾਇਆ ਨਹੀਂ ਗਿਆ ਤੇ \VUgras{ਲੱਕੜ ਦਾ ਸੰਦੂਕ}
\textsuperscript{(52)} \VUgras{ਬਾਲਣ} \textsuperscript{(53)} ਨਾਲ ਭਰਿਆ ਹੈ।

%%K t06-noto-1 noto
\VUnotes{143.2mm}{%
(*) Babo = ਨਿੱਕਾ ਬਾਲ; ਅੰਗਰੇਜ਼ੀ \textit{baby}, ਫ਼ਰਾਂਸੀਸੀ \textit{bébé},
ਇਤਾਲਵੀ \textit{bambino}।%
}

%%K t06-noto-2 noto
\VUnotes{143.2mm}{%
(*) Lambrequino = ਫ਼ਰਾਂਸੀਸੀ \textit{lambrequin}, ਸਪੇਨੀ
\textit{lambrequines}, ਪੁਰਤਗਾਲੀ \textit{lambrequins}, ਤੇ ਜਰਮਨ ਤੇ ਅੰਗਰੇਜ਼ੀ ਵਿਚ
ਵੀ \textit{lambrequins} --- ਅਰਥਾਤ ਜਰਮਨ, ਅੰਗਰੇਜ਼ੀ, ਫ਼ਰਾਂਸੀਸੀ, ਪੁਰਤਗਾਲੀ ਤੇ
ਸਪੇਨੀ, ਸਭਨਾਂ ਵਿਚ ਇਕੋ।%
}

%%K t06-c1-11-1 p
11. --- ਉਸ ਨੂੰ ਹਾਲੇ \VUgras{ਅੰਗੀਠੀ ਦੀ ਘੜੀ} \textsuperscript{(57)},
\VUgras{ਸ਼ਮ੍ਹਾਦਾਨ} \textsuperscript{(58)} ਤੇ \VUgras{ਸਜਾਵਟ ਦੀਆਂ ਤਸ਼ਤਰੀਆਂ}
\textsuperscript{(59)} ਝਾੜਨੀਆਂ ਪੈਣਗੀਆਂ, ਫਿਰ \VUgras{ਸ਼ੀਸ਼ਾ} \textsuperscript{(60)}
ਪੂੰਝਣਾ ਪਵੇਗਾ।

%%K t06-c1-12-1 p
12. --- ਰਿਹਾ ਬਾਲ ਦਾ \VUgras{ਪੰਘੂੜਾ} \textsuperscript{(61)} (ਉਹ babo (*)), ਉਹ ਤਾਂ
ਪਹਿਲਾਂ ਤੋਂ ਤਿਆਰ ਹੈ।

%%K t06-tit-3 sub
\VUsaut{5.0mm}
\VUpk{10.2pt}{40}{ਬੈਠਕ।}
\VUsaut{3.0mm}

%%K t06-c2-01-1 p
1. --- ਹੁਣ ਇਹ ਰਹੀ ਬੈਠਕ। ਇਹ ਬਹੁਤ ਸੋਹਣਾ \VUgras{ਕਮਰਾ} ਹੈ। ਸੱਜੇ ਕੋਨੇ ਦੀ ਅੰਗੀਠੀ ਇਕ
ਨਾਜ਼ੁਕ \VUgras{ਮੂਰਤੀ} \textsuperscript{(1)} ਤੇ ਦੋ ਸੋਹਣੇ \VUgras{ਫੁੱਲਦਾਨਾਂ}
\textsuperscript{(3)} ਨਾਲ ਸਜੀ ਹੈ, ਤੇ ਮਖ਼ਮਲ ਦੀ \VUgras{ਝਾਲਰ} \textsuperscript{(2)}
(*) ਨਾਲ ਢਕੀ ਹੈ।

%%K t06-c2-02-1 p
2. --- ਅੰਗੀਠੀ ਦੀਆਂ ਚਿੰਗਾੜੀਆਂ ਗ਼ਲੀਚੇ ਨੂੰ ਅੱਗ ਨਾ ਲਾ ਦੇਣ, ਇਸੇ ਲਈ ਉਸ ਦੇ ਸਾਹਮਣੇ
ਤਾਂਬੇ ਦਾ ਇਕ \VUgras{ਅੱਗ ਦਾ ਜੰਗਲਾ} \textsuperscript{(4)} ਰੱਖਿਆ ਗਿਆ ਹੈ।

%%K t06-c2-03-1 p
3. --- ਉਸ ਕੋਲ ਇਕ ਸਜੀਲੀ \VUgras{ਲਿਖਣ ਦੀ ਮੇਜ਼} \textsuperscript{(5)} ਖੁੱਲ੍ਹੀ ਖੜੀ ਹੈ।
ਇੱਥੇ ਹੀ \VUgras{ਘਰ ਦੀ ਮਾਲਕਣ} \textsuperscript{(23)} ਆਪਣੀਆਂ ਚਿੱਠੀਆਂ ਲਿਖਦੀ ਹੈ।

%%K t06-c2-04-1 p
4. --- ਖਿੜਕੀ ਉੱਤੇ \VUgras{ਜਾਲੀਦਾਰ ਪਰਦੇ} \textsuperscript{(6)} ਲੱਗੇ ਹਨ, ਜਿਨ੍ਹਾਂ
ਨੂੰ \VUgras{ਡੋਰੀਆਂ} \textsuperscript{(8)} ਬੰਨ੍ਹੀ ਰੱਖਦੀਆਂ ਹਨ, ਤੇ ਉਹ ਡੋਰੀਆਂ
\VUgras{ਖੂੰਟੀਆਂ} \textsuperscript{(7)} ਵਿਚ ਅਟਕੀਆਂ ਹਨ, ਤੇ ਉਨ੍ਹਾਂ ਉੱਤੇ ਭਾਰੇ ਕੱਪੜੇ
ਦੇ ਪਰਦੇ, ਜਿਨ੍ਹਾਂ ਦੇ ਉੱਤੇ ਇਕ \VUgras{ਝਾਲਰ-ਪੱਟੀ} \textsuperscript{(9)} ਹੈ।

%%K t06-c2-05-1 p
5. --- ਖਿੜਕੀ ਦੇ ਤਾਕ ਵਿਚ ਇਕ \VUgras{ਗਮਲਾ-ਚੌਕੀ} \textsuperscript{(11)} ਵਿਚ ਹਰਾ
\VUgras{ਬੂਟਾ} \textsuperscript{(10)} ਹੈ, ਇਕ ਸ਼ਾਨਦਾਰ ਖਜੂਰ, ਜਿਸ ਦੇ ਪੱਤੇ ਤਕਰੀਬਨ
\VUgras{ਤੇਲ ਦੇ ਲੈਂਪ} \textsuperscript{(12)} ਤਕ ਪਹੁੰਚਦੇ ਹਨ।

%%K t06-c2-06-1 p
6. --- \VUgras{ਲੈਂਪ ਦੀ ਛਾਂ} \textsuperscript{(13)} ਮਾਰੀਆ ਨੇ ਬਣਾਈ ਹੈ, ਉਹੋ ਮੋਹ ਲੈਣ
ਵਾਲੀ \VUgras{ਮੁਟਿਆਰ} \textsuperscript{(14)} ਜੋ \VUgras{ਪਿਆਨੋ} \textsuperscript{(15)}
ਉੱਤੇ ਬੈਠੀ ਹੈ। \VUgras{ਚੌਕੀ} \textsuperscript{(16)} ਉੱਤੇ ਸਿੱਧੀ ਬੈਠੀ ਉਹ ਇਕ ਧੁਨ ਦਾ
ਅਭਿਆਸ ਕਰ ਰਹੀ ਹੈ। ਉਹ ਬਹੁਤ ਸਿਆਣੀ ਤੇ ਬਹੁਤ ਸਲੀਕੇ ਵਾਲੀ ਹੈ, ਜਿਵੇਂ ਉਸ ਦੀ
\VUgras{ਸੁਰ-ਲਿਪੀ ਦੀ ਅਲਮਾਰੀ} \textsuperscript{(17)} ਤੋਂ ਪਤਾ ਲਗਦਾ ਹੈ: ਉਸ ਵਿਚ ਪੂਰੀ
ਤਰਤੀਬ ਹੈ।

%%K t06-tit-4 sub
\VUsaut{5.0mm}
\VUpk{10.2pt}{40}{ਸ਼ਰਾਰਤੀ।}
\VUsaut{3.0mm}

%%K t06-c2-07-1 p
7. --- ਉਸ ਦਾ ਭਰਾ ਪਾਲ ਇਕ \VUgras{ਕਿਤਾਬ} \textsuperscript{(19)} ਲੱਭ ਰਿਹਾ ਹੈ,
\VUgras{ਕਿਤਾਬਾਂ ਦੀ ਅਲਮਾਰੀ} \textsuperscript{(18)} ਵਿਚ। ਉਹ ਇਹੋ ਜਿਹਾ \VUgras{ਮੁੰਡਾ}
\textsuperscript{(20)} ਹੈ ਜੋ ਸ਼ਰਾਰਤੀ ਤੇ ਬਿਲਕੁਲ ਬੇਕਾਬੂ ਹੈ। ਬਹੁਤ ਉੱਚੀ ਰੱਖੀ ਕਿਤਾਬ ਤਕ
ਪਹੁੰਚਣ ਲਈ ਅਲਮਾਰੀ ਤੋਂ ਲਟਕਦਾ ਹੋਇਆ ਉਹ ਉਸ ਦੇ ਉੱਤੇ ਰੱਖੀ \VUgras{ਮੂਰਤੀ}
\textsuperscript{(21)} ਡੇਗ ਦੇਣ ਦਾ ਖ਼ਤਰਾ ਚੁੱਕ ਰਿਹਾ ਹੈ।

%%K t06-c2-08-1 p
8. --- ਇਹ ਸ਼ਰਾਰਤ ਕਰਨ ਲਈ ਉਹ ਇਸ ਗੱਲ ਦਾ ਫ਼ਾਇਦਾ ਚੁੱਕ ਰਿਹਾ ਹੈ ਕਿ ਉਸ ਦੀ ਮਾਂ ਇਕ ਸਹੇਲੀ
ਨਾਲ ਗੱਲਾਂ ਵਿਚ ਲੱਗੀ ਹੈ, ਇਕ \VUgras{ਬੀਬੀ} \textsuperscript{(22)} ਜੋ ਕੁਝ ਦਿਨ ਉਨ੍ਹਾਂ
ਨਾਲ ਲੰਘਾਉਣ ਆਈ ਹੈ। ਮਾਂ \VUgras{ਸੋਫ਼ੇ} \textsuperscript{(24)} ਉੱਤੇ ਬੈਠੀ ਹੈ,
\VUgras{ਆਰਾਮ-ਕੁਰਸੀ} \textsuperscript{(25)} ਕੋਲ, ਤੇ ਉਸ ਦੀ ਸਹੇਲੀ ਉਸ \VUgras{ਕੁਰਸੀ}
\textsuperscript{(27)} ਉੱਤੇ, ਜਿਸ ਦੀ ਸਾਨੂੰ ਸਿਰਫ਼ \VUgras{ਪਿੱਠ}
\textsuperscript{(26)} ਵਿਖਾਈ ਦਿੰਦੀ ਹੈ।

%%K t06-c2-09-1 p
9. --- ਕੰਧ ਉੱਤੇ ਟੰਗੇ ਵੱਡੇ \VUgras{ਫ਼ਰੇਮ} \textsuperscript{(30)} ਵਿਚ ਇਕ ਰਿਸ਼ਤੇਦਾਰ
ਦੀ \VUgras{ਤਸਵੀਰ} \textsuperscript{(29)} ਹੈ। ਮੇਜ਼ ਉੱਤੇ ਪਈ \VUgras{ਐਲਬਮ}
\textsuperscript{(32)} ਵਿਚ ਬਾਕੀ ਪਰਿਵਾਰ ਦੀਆਂ \VUgras{ਤਸਵੀਰਾਂ}
\textsuperscript{(33)} ਇਕੱਠੀਆਂ ਹਨ। ਬੱਚੇ ਹੁਣੇ ਥੋੜੀ ਦੇਰ ਪਹਿਲਾਂ ਉਸ ਨੂੰ ਪੁੱਠ-ਪਰਤ ਰਹੇ
ਸਨ, ਕਿਉਂਕਿ \VUgras{ਕੁਰਸੀਆਂ} \textsuperscript{(31)} ਹਾਲੇ ਵੀ ਮੇਜ਼ ਦੇ ਚੁਫੇਰੇ ਜਮ੍ਹਾਂ
ਹਨ।

%%K t06-c2-10-1 p
10. --- ਚੀਨੀ ਮਿੱਟੀ ਦਾ \VUgras{ਚੀਨੀ ਫੁੱਲਦਾਨ} \textsuperscript{(34)}, ਜੋ
\VUgras{ਕਾਗਜ਼ ਕੱਟਣ ਦੀ ਛੁਰੀ} \textsuperscript{(35)} ਕੋਲ ਰੱਖਿਆ ਹੈ, ਮਾਰੀਆ ਨੂੰ ਉਸ ਦੇ
ਨਾਂ ਦੇ ਦਿਨ ਮਿਲਿਆ ਸੀ, ਤੇ ਕਮਰੇ ਦਾ ਕੋਨਾ ਭਰਦਾ \VUgras{ਤਹਿ ਹੋਣ ਵਾਲਾ ਪਰਦਾ}
\textsuperscript{(36)} ਉਸ ਦੇ ਚਾਚਾ ਚੀਨ ਤੋਂ ਲਿਆਏ ਸਨ।

%%K t06-c2-11-1 p
11. --- ਸੋਹਣੀਆਂ \VUgras{ਤਸਵੀਰਾਂ} \textsuperscript{(37)} ਨਾਲ ਜਗਮਗਾਉਂਦੀ, ਜੋ
\VUgras{ਕੰਧ} \textsuperscript{(42)} ਉੱਤੇ ਟੰਗੀਆਂ ਹਨ, ਛੱਤ ਦੇ \VUgras{ਝਾੜ-ਫ਼ਾਨੂਸ}
\textsuperscript{(38)} ਨਾਲ ਰੌਸ਼ਨ, ਜਿਸ ਦੇ \VUgras{ਬਿਜਲੀ ਦੇ ਲੈਂਪ}
\textsuperscript{(39)} ਸ਼ਾਮ ਨੂੰ ਇੰਨੇ ਖ਼ੁਸ਼ ਚਮਕਦੇ ਹਨ, ਇਹ ਬੈਠਕ ਬੜੀ ਸੋਹਣੀ ਲਗਦੀ ਹੈ।
ਫ਼ਰਸ਼ ਉੱਤੇ ਵਿਛਿਆ ਕੋਮਲ \VUgras{ਗ਼ਲੀਚਾ} \textsuperscript{(40)} ਤੇ ਇੰਨੀ ਕੰਮ ਦੀ
\VUgras{ਬਿਜਲੀ ਦੀ ਘੰਟੀ} \textsuperscript{(41)} ਉਸ ਨੂੰ ਪੂਰੀ ਤਰ੍ਹਾਂ ਆਰਾਮਦੇਹ ਬਣਾ ਦਿੰਦੇ
ਹਨ।

%%K t06-tit-5 sub
\VUsaut{3.6mm}
\VUpk{10.2pt}{40}{ਖਾਣੇ ਦਾ ਕਮਰਾ।}
\VUsaut{3.0mm}

%%K t06-c3-01-1 p
1. --- ਖਾਣੇ ਦੀ ਘੰਟੀ ਵੱਜ ਚੁੱਕੀ ਹੈ। ਮਾਰੀਆ ਨੂੰ ਛੱਡ ਕੇ ਸਾਰਾ ਪਰਿਵਾਰ ਮੇਜ਼ ਦੇ ਚੁਫੇਰੇ
ਜਮ੍ਹਾਂ ਹੈ। ਖਾਣੇ ਦਾ ਕਮਰਾ ਇਕ ਵਧੀਆ ਮਿੱਟੀ ਦੇ \VUgras{ਚੁੱਲ੍ਹੇ} \textsuperscript{(1)}
ਨਾਲ ਸੋਹਣਾ ਨਿੱਘਾ ਹੈ, ਜਿਸ ਦੀ \VUgras{ਨਾਲੀ} \textsuperscript{(2)} ਪਤਲੀ ਹੈ।

%%K t06-c3-02-1 p
2. --- ਇਸ ਕਮਰੇ ਵਿਚ ਸਾਮਾਨ ਬਹੁਤ ਨਹੀਂ। ਕੰਧ ਨਾਲ, \VUgras{ਤਿਪਾਈ}
\textsuperscript{(4)} ਦੇ ਉੱਤੇ, ਇਕ ਸਜਾਵਟੀ ਨਿੱਕਾ \VUgras{ਫੁਹਾਰਾ}
\textsuperscript{(3)} ਟੰਗਿਆ ਹੈ। ਕੋਲ ਹੀ \VUgras{ਸਜਾਵਟੀ ਅਲਮਾਰੀ} \textsuperscript{(5)}
ਹੈ, ਜਿਸ ਉੱਤੇ ਠੰਢੇ ਖਾਣੇ, ਮਿਠਾਈਆਂ ਤੇ ਉਹ ਭਾਂਡੇ ਰੱਖੇ ਜਾਂਦੇ ਹਨ ਜਿਨ੍ਹਾਂ ਦੀ ਇਸ ਵੇਲੇ
ਲੋੜ ਨਹੀਂ।

%%K t06-c3-03-1 p
3. --- ਸੋ ਉੱਪਰਲੇ ਖ਼ਾਨੇ ਵਿਚ ਸਾਨੂੰ ਇਕ \VUgras{ਭੁੰਨਿਆ ਕੁੱਕੜ} \textsuperscript{(7)}
ਵਿਖਾਈ ਦਿੰਦਾ ਹੈ, ਇਕ \VUgras{ਮੱਛੀ} \textsuperscript{(8)} ਤੇ ਇਕ \VUgras{ਕਟੋਰਾ}
\textsuperscript{(10)} \VUgras{ਫਲਾਂ} \textsuperscript{(9)} ਦਾ। ਹੇਠਾਂ \VUgras{ਪਨੀਰ}
\textsuperscript{(11)}, \VUgras{ਰੋਟੀ} \textsuperscript{(12)} ਤੇ
\VUgras{ਤੇਲ-ਸਿਰਕੇ ਦਾ ਦਾਨ} \textsuperscript{(13)} ਹੈ, ਜਿਸ ਵਿਚ ਸਲਾਦ ਸਜਾਉਣ ਦਾ ਸਭ ਕੁਝ
ਰਹਿੰਦਾ ਹੈ: ਤੇਲ, ਸਿਰਕਾ, \VUgras{ਲੂਣ} \textsuperscript{(14)} ਤੇ \VUgras{ਮਿਰਚ}
\textsuperscript{(15)}। ਅੰਤ ਵਿਚ ਸਭ ਤੋਂ ਹੇਠਲੇ ਖ਼ਾਨੇ ਵਿਚ ਇਕ \VUgras{ਕੇਕ}
\textsuperscript{(17)} ਹੈ, \VUgras{ਰਾਈ ਦੇ ਦਾਨ} \textsuperscript{(16)} ਕੋਲ।

%%K t06-c3-04-1 p
4. --- ਖਾਣੇ ਦੇ ਕਮਰੇ ਦੀ ਘੜੀ, ਜਿਸ ਨੂੰ \VUgras{ਕੰਧ-ਘੜੀ} \textsuperscript{(20)}
ਆਖਦੇ ਹਨ, ਬਹੁਤ ਅਨੋਖੀ ਸ਼ਕਲ ਦੀ ਹੈ। ਉਹ ਲੱਕੜ ਦੇ ਇਕ ਚੌਖਟੇ ਵਿਚ ਜੜੀ ਹੈ, ਜਿਸ ਵਿਚ ਇਕ
\VUgras{ਹਵਾ-ਦਬਾਅ-ਮਾਪ} \textsuperscript{(19)} ਤੇ ਇਕ \VUgras{ਤਾਪ-ਮਾਪ}
\textsuperscript{(18)} ਵੀ ਹਨ।

%%K t06-tit-6 sub
\VUsaut{4.6mm}
\VUpk{10.2pt}{40}{ਖਾਣਾ ਵਰਤਾਉਣਾ।}
\VUsaut{3.0mm}

%%K t06-c3-05-1 p
5. --- ਕਮਰੇ ਦੇ ਚੁਫੇਰੇ ਕੰਧਾਂ ਉੱਤੇ ਮੋਮ ਦੀ ਪਾਲਿਸ਼ ਕੀਤੀ ਬਲੂਤ ਦੀ
\VUgras{ਲੱਕੜ ਦੀ ਚੌਖਟ} \textsuperscript{(21)} ਲੱਗੀ ਹੈ। ਕੱਪੜੇ ਦਾ ਇਕ
\VUgras{ਬੂਹੇ ਦਾ ਪਰਦਾ} \textsuperscript{(22)} ਬਰਾਂਡੇ ਵੱਲ ਜਾਂਦੇ ਬੂਹੇ ਨੂੰ ਚੰਗੀ
ਤਰ੍ਹਾਂ ਢਕ ਦਿੰਦਾ ਹੈ। ਖਾਣੇ ਪਿੱਛੋਂ ਪਰਿਵਾਰ ਉੱਥੇ ਹੀ ਕੌਫ਼ੀ ਪੀਣ ਜਾਵੇਗਾ।
\VUgras{ਪਿਆਲੇ} \textsuperscript{(24)} ਤੇ \VUgras{ਤਸ਼ਤਰੀਆਂ} \textsuperscript{(25)}
ਪਹਿਲਾਂ ਤੋਂ \VUgras{ਚੀਨੀ ਭਾਂਡਿਆਂ ਦੀ ਅਲਮਾਰੀ} \textsuperscript{(23)} ਉੱਤੇ ਸਜੇ ਹਨ।

%%K t06-c3-06-1 p
6. --- ਮੇਜ਼ ਦੇ ਉੱਤੇ ਛੱਤ ਤੋਂ ਇਕ \VUgras{ਲਟਕਦਾ ਲੈਂਪ} \textsuperscript{(26)}
ਬੰਨ੍ਹਿਆ ਹੈ; ਉਸ ਦੀ ਬਹੁਤ ਤੇਜ਼ ਲੋਅ ਸ਼ਾਮ ਨੂੰ ਸਾਰੇ ਖਾਣੇ ਦੇ ਕਮਰੇ ਵਿਚ ਭਰ ਜਾਂਦੀ ਹੈ।

%%K t06-c3-07-1 p
7. --- ਹੁਣ ਆਪ \VUgras{ਖਾਣੇ ਦੀ ਮੇਜ਼} \textsuperscript{(27)} ਵੇਖੀਏ। ਪਰਿਵਾਰ ਦੇ ਆਉਣ
ਉੱਤੇ ਮੇਜ਼ ਲੱਗੀ ਹੋਈ ਸੀ। \VUgras{ਮੇਜ਼ਪੋਸ਼} \textsuperscript{(28)} ਉੱਤੇ ਹਰ ਖਾਣ ਵਾਲੇ
ਦੀ ਥਾਂ ਰੱਖੀ ਸੀ ਇਕ \VUgras{ਪਲੇਟ} \textsuperscript{(33)}, ਇਕ \VUgras{ਰੁਮਾਲ}
\textsuperscript{(29)}, ਇਕ \VUgras{ਚਮਚਾ} \textsuperscript{(34)}, ਇਕ \VUgras{ਕਾਂਟਾ}
\textsuperscript{(35)}, ਇਕ \VUgras{ਛੁਰੀ} \textsuperscript{(36)} ਤੇ ਇਕ
\VUgras{ਗਿਲਾਸ} \textsuperscript{(32)}। ਸ਼ਰਾਬ ਦੀਆਂ \VUgras{ਬੋਤਲਾਂ}
\textsuperscript{(30)} ਤੇ ਪਾਣੀ ਦੀ ਇਕ \VUgras{ਸੁਰਾਹੀ} \textsuperscript{(31)} ਵੀ ਸੀ।

%%K t06-c3-08-1 p
8. --- \VUgras{ਨੌਕਰ} \textsuperscript{(39)} ਪਹਿਲਾਂ \VUgras{ਸ਼ੋਰਬੇ ਦਾ ਭਾਂਡਾ}
\textsuperscript{(37)} ਲਿਆਇਆ, ਜਿਸ ਵਿਚੋਂ ਹਾਲੇ ਵੀ ਕੁਝ ਸ਼ੋਰਬੇ ਤੋਂ ਭਾਫ਼ ਉੱਠ ਰਹੀ ਹੈ।
ਹੁਣ ਉਹ ਪਹਿਲਾ \VUgras{ਖਾਣਾ} \textsuperscript{(38)} ਵਰਤਾ ਰਿਹਾ ਹੈ, ਮਾਸ ਦਾ। ਫਿਰ ਉਹ ਉਹ
ਖਾਣੇ ਲਿਆਵੇਗਾ ਜਿਨ੍ਹਾਂ ਦੇ ਨਾਂ ਅਸੀਂ ਲੈ ਚੁੱਕੇ ਹਾਂ; ਤੇ ਅੰਤ ਵਿਚ, \VUgras{ਮਿੱਠੇ}
ਪਿੱਛੋਂ, ਉਹ ਇਸ ਗੱਲ ਦਾ ਫ਼ਾਇਦਾ ਚੁੱਕੇਗਾ ਕਿ ਮਾਲਕ ਬਰਾਂਡੇ ਵਿਚ ਚਲੇ ਗਏ ਹਨ, ਤੇ ਮੇਜ਼ ਖ਼ਾਲੀ
ਕਰ ਕੇ ਖਾਣੇ ਦਾ ਕਮਰਾ ਮੁੜ ਠੀਕ ਕਰ ਦੇਵੇਗਾ।

%%K t06-c3-08-2 p
\textit{ਟਿੱਪਣੀ।} --- \textit{ਖਾਣੇ ਨਾਲ ਜੁੜੇ ਰੋਜ਼ ਦੇ ਸ਼ਬਦ ਇੱਥੇ ਸਿਰਫ਼ ਮੋਟੇ ਤੌਰ ਉੱਤੇ
ਦਿੱਤੇ ਗਏ ਹਨ। ਇਹ ਸੂਚੀ « ਹੋਟਲ » (ਨੰ. 13) ਤੇ « ਮੰਡੀ » (ਨੰ. 15) ਦੀਆਂ ਦੋ ਸਾਰਣੀਆਂ ਵਿਚ
ਪੂਰੀ ਕੀਤੀ ਜਾਵੇਗੀ।}

%%K t06-tit-7 sub
\VUsaut{4.6mm}
\VUpk{10.2pt}{40}{ਰਸੋਈ।}
\VUsaut{3.0mm}

%%K t06-c4-01-1 p
1. --- ਅੱਧ-ਖੁੱਲ੍ਹੇ ਬੂਹੇ ਵਿਚੋਂ ਜਿਸ ਰਸੋਈ ਉੱਤੇ ਅਸੀਂ ਨਜ਼ਰ ਮਾਰਦੇ ਹਾਂ, ਉੱਥੇ ਇਕ ਸੋਹਣਾ
ਗੜਬੜ-ਝਾਲਾ ਲੱਭਦਾ ਹੈ। \VUgras{ਚੁੱਲ੍ਹੇ} \textsuperscript{(1)} ਦੇ ਸਾਹਮਣੇ, ਜਿੱਥੇ
\VUgras{ਅੱਗ} \textsuperscript{(3)} ਚਟਕ ਰਹੀ ਹੈ, \VUgras{ਭੁੰਨਣ ਦਾ ਜੰਤਰ}
\textsuperscript{(5)} ਲੱਗਾ ਹੈ। ਇਕ ਵਧੀਆ \VUgras{ਭੁੰਨਿਆ ਮਾਸ} \textsuperscript{(7)}
\VUgras{ਸੀਖ਼} \textsuperscript{(6)} ਉੱਤੇ ਹੈ ਤੇ ਪੱਕਣ ਨੂੰ ਹੈ।

%%K t06-c4-02-1 p
2. --- \VUgras{ਧੌਂਕਣੀ} \textsuperscript{(8)} ਤੇ \VUgras{ਕੋਲੇ ਦੀ ਕਹੀ}
\textsuperscript{(9)}, ਜੋ ਹੁਣੇ ਕੰਮ ਆਈਆਂ ਹਨ, \VUgras{ਲੱਕੜਾਂ} \textsuperscript{(10)}
ਸਮੇਤ ਭੁੰਜੇ ਹੀ ਛੱਡ ਦਿੱਤੀਆਂ ਗਈਆਂ ਹਨ।

%%K t06-c4-03-1 p
3. --- ਅੰਗੀਠੀ ਦੇ ਉੱਤੇ ਦਾ ਹਿੱਸਾ ਸੰਵਰਿਆ ਹੋਇਆ ਹੈ: \VUgras{ਲਾਲਟੈਣ}
\textsuperscript{(11)}, \VUgras{ਦੀਆਸਲਾਈ-ਦਾਨ} \textsuperscript{(13)} ਜਿਸ ਵਿਚ
\VUgras{ਦੀਆਸਲਾਈਆਂ} \textsuperscript{(12)} ਹਨ, \VUgras{ਸ਼ਮ੍ਹਾਦਾਨ}
\textsuperscript{(14)} ਤੇ \VUgras{ਹੱਥ ਦਾ ਸ਼ਮ੍ਹਾਦਾਨ} \textsuperscript{(15)} ਆਪਣੀ
\VUgras{ਮੋਮਬੱਤੀ} \textsuperscript{(16)} ਸਮੇਤ — ਸਭ ਆਪਣੀ ਥਾਂ ਹਨ। ਪਰ ਵੱਡੀ
\VUgras{ਕੋਲੇ ਦੀ ਟੋਕਰੀ} \textsuperscript{(18)}, ਜੋ ਉਸ \VUgras{ਕੋਲੇ}
\textsuperscript{(17)} ਨਾਲ ਭਰੀ ਹੈ ਜੋ \VUgras{ਰਸੋਇਣ} \textsuperscript{(23)} ਹੁਣੇ
ਤਹਿਖ਼ਾਨੇ ਤੋਂ ਲਿਆਈ ਹੈ, ਰਸੋਈ ਦੇ ਵਿਚਕਾਰ ਹੀ ਪਈ ਰਹਿ ਗਈ ਹੈ।

%%K t06-c4-04-1 p
4. --- ਸੱਚ ਤਾਂ ਇਹ ਹੈ ਕਿ ਵਿਚਾਰੀ ਨੂੰ ਛੇਤੀ ਕਰਨੀ ਹੈ, ਜੇ ਦੁਪਹਿਰ ਦਾ ਖਾਣਾ ਵੇਲੇ ਸਿਰ
ਤਿਆਰ ਕਰਨਾ ਹੈ। ਇਸ ਵੇਲੇ ਉਹ \VUgras{ਪਕਾਉਣ ਦੇ ਚੁੱਲ੍ਹੇ} \textsuperscript{(19)} ਉੱਤੇ ਇਕ
\VUgras{ਤਰੀ} \textsuperscript{(24)} ਚਲਾ ਰਹੀ ਹੈ, ਜੋ ਸੜਨ ਨਾ ਪਾਵੇ।

%%K t06-c4-05-1 p
5. --- \VUgras{ਤੰਦੂਰ} \textsuperscript{(20)} ਵਿਚ ਇਕ ਕੇਕ ਪੱਕ ਰਿਹਾ ਹੈ, ਜੋ ਉਸ ਨੇ
ਬੱਚਿਆਂ ਦੀ ਸ਼ਾਮ ਦੀ ਚਾਹ ਲਈ ਬਣਾਇਆ ਹੈ। ਚੁੱਲ੍ਹੇ ਦੇ ਉੱਤੇ \VUgras{ਕੜਛੀ}
\textsuperscript{(21)} ਤੇ \VUgras{ਝਰਨਾ} \textsuperscript{(22)} ਟੰਗੇ ਹਨ।

%%K t06-tit-8 sub
\VUsaut{4.0mm}
\VUpk{10.2pt}{40}{ਰਸੋਈ ਦੇ ਭਾਂਡੇ।}
\VUsaut{3.0mm}

%%K t06-c4-06-1 p
6. --- ਕਮਰੇ ਦੇ ਪਿੱਛੇ ਟੰਗਿਆ \VUgras{ਭਾਂਡਿਆਂ ਦਾ ਸੈੱਟ} \textsuperscript{(25)} ਬਹੁਤ
ਸਾਫ਼ ਹੈ। ਤਾਂਬੇ ਦੀਆਂ ਸੋਹਣੀਆਂ \VUgras{ਪਤੀਲੀਆਂ} \textsuperscript{(26)},
\VUgras{ਦੁੱਧ ਦਾ ਜੱਗ} \textsuperscript{(27)}, \VUgras{ਛਾਨਣੀ} \textsuperscript{(28)}
ਤੇ \VUgras{ਕੱਦੂਕਸ} \textsuperscript{(29)} ਧਿਆਨ ਨਾਲ ਮਾਂਜੇ ਗਏ ਹਨ।

%%K t06-c4-07-1 p
7. --- \VUgras{ਲੂਣ-ਦਾਨ} \textsuperscript{(30)} ਨਿੱਕੀ ਅਲਮਾਰੀ ਉੱਤੇ \VUgras{ਚਾਹ-ਦਾਨ}
\textsuperscript{(31)} ਤੇ \VUgras{ਤੇਲ ਦੀ ਵੱਡੀ ਬੋਤਲ} \textsuperscript{(32)} ਕੋਲ ਰੱਖਿਆ
ਹੈ। \VUgras{ਦਰਾਜ਼} \textsuperscript{(33)} ਵਿਚ ਸ਼ਾਇਦ ਚਮਚੇ-ਕਾਂਟੇ ਤੇ ਰਸੋਈ ਦੀਆਂ ਛੁਰੀਆਂ
ਰਹਿੰਦੀਆਂ ਹਨ।

%%K t06-c4-08-1 p
8. --- \VUgras{ਝਾੜੂ} \textsuperscript{(34)} ਤੇ \VUgras{ਖੰਭਾਂ ਦਾ ਝਾੜਨ}
\textsuperscript{(35)} ਇਕ ਕੋਨੇ ਵਿਚ ਹਨ। ਰਸੋਇਣ ਹੁਣੇ \VUgras{ਕੱਟਣ ਦਾ ਗੁੱਡਾ}
\textsuperscript{(36)} ਕੰਮ ਵਿਚ ਲਿਆ ਚੁੱਕੀ ਹੈ; ਉਸ ਨੇ ਉਸ ਨੂੰ ਖਿੜਕੀ ਕੋਲ ਛੱਡ ਦਿੱਤਾ ਹੈ।

%%K t06-c4-09-1 p
9. --- ਇਕ \VUgras{ਹੱਥ ਦਾ ਤੌਲੀਆ} \textsuperscript{(37)} ਕੰਧ ਉੱਤੇ ਟੰਗਿਆ ਹੈ,
\VUgras{ਕੀਪ} \textsuperscript{(38)} ਕੋਲ। ਰਸੋਈ ਦੇ ਇਸ ਹਿੱਸੇ ਵਿਚ \VUgras{ਸੀਖ਼-ਜਾਲੀ}
\textsuperscript{(39)}, \VUgras{ਕਟਾਰ} \textsuperscript{(40)} ਤੇ
\VUgras{ਕੱਟਣ ਦਾ ਪਟੜਾ} \textsuperscript{(41)} ਰੱਖੇ ਜਾਂਦੇ ਹਨ। ਫਿਰ ਕਟਾਰ ਦੇ ਉੱਤੇ, ਇਕ
\VUgras{ਤਾਕ} \textsuperscript{(42)} ਉੱਤੇ, \VUgras{ਹਾਂਡੀਆਂ} \textsuperscript{(43)},
\VUgras{ਸਾਲਣ ਦੀ ਦੇਗਚੀ} \textsuperscript{(44)} ਆਪਣੇ \VUgras{ਢੱਕਣ}
\textsuperscript{(45)} ਸਮੇਤ, ਤੇ \VUgras{ਕੜਾਹਾ} \textsuperscript{(46)} ਹਨ।
\VUgras{ਤਲਣ ਦੀ ਕੜਾਹੀ} \textsuperscript{(47)} ਕੋਲ ਹੀ ਟੰਗੀ ਹੈ।

%%K t06-c4-10-1 p
10. --- ਇਸ ਕੋਨੇ ਵਿਚ ਸਿਰਫ਼ ਤਾਂਬੇ ਦੀ \VUgras{ਟੂਟੀ} \textsuperscript{(48)} ਚਮਕ
ਪਾਉਂਦੀ ਹੈ। \VUgras{ਨਾਲੀ-ਚੌਕੀ} \textsuperscript{(49)}, ਜਿਸ ਉੱਤੇ ਵੱਡਾ \VUgras{ਘੜਾ}
\textsuperscript{(50)} ਰੱਖਿਆ ਹੈ, ਆਪਣੀ ਨਿੱਕੀ ਅਲਮਾਰੀ ਨਾਲ ਬਹੁਤ ਕੰਮ ਦੀ ਹੋਵੇਗੀ, ਜਿੱਥੇ
\VUgras{ਸਿਰਕੇ ਦਾ ਪੀਪਾ} \textsuperscript{(51)} ਆਪਣੀ \VUgras{ਟੂਟੀ}
\textsuperscript{(52)} ਸਮੇਤ, \VUgras{ਕੂੜੇ ਦਾ ਸੰਦੂਕ} \textsuperscript{(53)} ਤੇ
\VUgras{ਗੰਦੇ ਭਾਂਡੇ} \textsuperscript{(54)} ਰੱਖੇ ਜਾਂਦੇ ਹਨ।

%%K t06-tit-9 sub
\VUsaut{4.0mm}
\VUpk{10.2pt}{40}{ਰਸੋਈ ਵਿਚ ਰੱਖੀਆਂ ਹੋਰ ਚੀਜ਼ਾਂ।}
\VUsaut{3.0mm}

%%K t06-c4-11-1 p
11. --- ਉਹ ਵੱਡਾ \VUgras{ਟੱਬ} \textsuperscript{(55)} ਜਿਸ ਵਿਚ \VUgras{ਸਬਜ਼ੀਆਂ}
\textsuperscript{(58)} ਧੋਤੀਆਂ ਜਾਂਦੀਆਂ ਹਨ, ਤੇ ਢਲੇ ਲੋਹੇ ਦਾ \VUgras{ਕੜਾਹਾ}
\textsuperscript{(56)}, ਜੋ \VUgras{ਇੱਟ ਦੇ ਫ਼ਰਸ਼} \textsuperscript{(57)} ਉੱਤੇ ਰੱਖਿਆ
ਹੈ, ਸ਼ਾਇਦ ਉਸੇ ਅਲਮਾਰੀ ਦੇ ਹਨ।

%%K t06-c4-12-1 p
12. --- ਰਸੋਇਣ ਕੋਲ ਚੁੱਲ੍ਹਿਆਂ ਦੀ ਘਾਟ ਨਹੀਂ, ਕਿਉਂਕਿ ਇਹ ਰਿਹਾ ਮੇਜ਼ ਦੇ ਕੋਨੇ ਉੱਤੇ ਇਕ
\VUgras{ਗੈਸ ਦਾ ਚੁੱਲ੍ਹਾ} \textsuperscript{(59)}, ਜਿਸ ਉੱਤੇ ਇਕ ਨਿੱਕੀ ਪਤੀਲੀ ਵਿਚ ਪਾਣੀ
ਗਰਮ ਹੋ ਰਿਹਾ ਹੈ। ਉਸ ਵਿਚੋਂ ਨਿਕਲਦੀ \VUgras{ਭਾਫ਼} \textsuperscript{(60)} ਦੱਸਦੀ ਹੈ ਕਿ
ਉਹ ਉਬਲ ਰਿਹਾ ਹੋਵੇਗਾ। ਇਹ ਪਾਣੀ ਸ਼ਾਇਦ ਕੌਫ਼ੀ ਲਈ ਹੈ, ਕਿਉਂਕਿ ਮੇਜ਼ ਦੇ ਦੂਜੇ ਸਿਰੇ ਉੱਤੇ ਇਹ
ਰਹੇ \VUgras{ਕੌਫ਼ੀ ਦਾ ਭਾਂਡਾ} \textsuperscript{(64)} ਤੇ \VUgras{ਕੌਫ਼ੀ ਦੀ ਚੱਕੀ}
\textsuperscript{(63)}।

%%K t06-c4-13-1 p
13. --- ਗੈਸ ਦਾ ਚੁੱਲ੍ਹਾ \VUgras{ਗੈਸ ਦੀ ਟੂਟੀ} \textsuperscript{(62)} ਨਾਲ ਜੁੜਿਆ ਹੈ,
ਇਕ ਲੰਮੀ \VUgras{ਰਬੜ ਦੀ ਨਾਲੀ} \textsuperscript{(61)} ਦੇ ਜ਼ਰੀਏ।

%%K t06-c4-14-1 p
14. --- \VUgras{ਦਿਹਾੜੀ ਦੀ ਸਹਾਇਕ} \textsuperscript{(66)}, ਜੋ ਰਸੋਇਣ ਦੀ ਮਦਦ ਨੂੰ
ਆਉਂਦੀ ਹੈ, ਰਾਤ ਦੇ ਖਾਣੇ ਦੀਆਂ ਸਬਜ਼ੀਆਂ ਤਿਆਰ ਕਰ ਰਹੀ ਹੈ। ਛਿੱਲ ਕੇ ਤੇ ਧੋ ਕੇ ਉਹ ਉਨ੍ਹਾਂ
ਨੂੰ \VUgras{ਭਾਂਡੇ} \textsuperscript{(65)} ਵਿਚ ਰੱਖ ਦੇਵੇਗੀ, ਜਦੋਂ ਤਕ ਪਕਾਉਣ ਦਾ ਵੇਲਾ ਨਾ
ਆਵੇ।

%%K t06-tit-10 sub
\VUsaut{4.0mm}
{\centering\fontsize{10.2pt}{10.2pt}\selectfont\textls[40]{\scshape ਗ਼ੁਸਲਖ਼ਾਨਾ।} \textit{(ਨ੍ਹਾਉਣ ਦਾ ਕਮਰਾ।)}\par}
\VUsaut{3.0mm}

%%K t06-c5-01-1 p
1. --- ਘਰ ਦੇ ਮੁੱਖ ਕਮਰੇ ਜਾਣ ਲੈਣ ਲਈ ਹੁਣ ਇਕੋ ਝਾਤੀ ਬਾਕੀ ਹੈ: ਗ਼ੁਸਲਖ਼ਾਨੇ ਦਾ ਸਾਮਾਨ
ਵੇਖਣਾ। ਇਸ ਵਿਚ ਦੇਰ ਨਾ ਲੱਗੇਗੀ, ਕਿਉਂਕਿ ਉਹ ਵੱਡਾ ਨਹੀਂ, ਤੇ ਅਸੀਂ ਝੱਟ ਵੇਖ ਲਵਾਂਗੇ ਕਿ
\VUgras{ਨ੍ਹਾਉਣ ਦੇ ਟੱਬ} \textsuperscript{(1)} ਨਾਲ ਇਕ ਚੰਗਾ \VUgras{ਪਾਣੀ ਗਰਮ ਕਰਨ ਦਾ
ਜੰਤਰ} \textsuperscript{(2)} ਹੈ, ਕਿ \VUgras{ਸਾਬਣਦਾਨੀ} \textsuperscript{(3)} ਨ੍ਹਾਉਣ
ਵਾਲੇ ਦੇ ਹੱਥ ਦੀ ਪਹੁੰਚ ਵਿਚ ਹੈ, ਤੇ ਕਿ \VUgras{ਤੌਲੀਏ ਦੇ ਡੰਡੇ} \textsuperscript{(5)}
ਨਾਲ \VUgras{ਤੌਲੀਏ} \textsuperscript{(4)} ਸੁਕਾਉਣਾ ਸੌਖਾ ਹੁੰਦਾ ਹੋਵੇਗਾ।

%%K t06-c5-02-1 p
2. --- ਇਸ ਕਮਰੇ ਦੀਆਂ ਕੰਧਾਂ \VUgras{ਚਮਕਦੀਆਂ ਟਾਈਲਾਂ} \textsuperscript{(6)} ਨਾਲ ਢਕੀਆਂ
ਹਨ, ਜਿਸ ਨਾਲ ਇਕ \VUgras{ਫੁਹਾਰਾ} \textsuperscript{(7)} ਲਾਉਣਾ ਸੰਭਵ ਹੋਇਆ, ਬਿਨਾਂ ਇਸ ਡਰ
ਦੇ ਕਿ ਉਸ ਨੂੰ ਵਰਤਦੇ ਵੇਲੇ ਛਿੱਟੇ ਉਨ੍ਹਾਂ ਨੂੰ ਵਿਗਾੜ ਦੇਣ। ਪੈਰ ਧੋਣ ਦਾ ਜਸਤ ਦਾ ਇਕ ਨਿੱਕਾ
\VUgras{ਤਸਲਾ} \textsuperscript{(8)} ਸਾਮਾਨ ਪੂਰਾ ਕਰ ਦਿੰਦਾ ਹੈ।
\VUsaut{6.0mm}
\VUfilet{22mm}
